\documentclass{galahad}

% set the package name

\newcommand{\package}{slblt}
\newcommand{\packagename}{SLBLT}
\newcommand{\fullpackagename}{\libraryname\_\-\packagename}

\begin{document}

\galheader

%%%%%%%%%%%%%%%%%%%%%% SUMMARY %%%%%%%%%%%%%%%%%%%%%%%%

\galsummary

This package 
{\bf solves a sparse symmetric system of linear equations.} Given 
a sparse symmetric matrix $\bmA = \{ a_{ij} \}_{n \times n}$, and an 
$n$-vector $\bmr$ or a matrix $\bmR = \{ r_{ij} \}_{n \times r}$
of right-hand sides, this package
solves the system $\bmA \bmx = \bmr$ or the system 
$\bmA \bmX = \bmR$  using a multi-frontal factorization method.
The matrix $\bmA$ need not be definite. Specifically,
the following cases are covered:

\noindent
1. $\bmA$ is {\bf indefinite}.
In this case, \packagename\ computes the sparse factorization
\[
\bmA =  \bmP\bmL\bmB(\bmP\bmL)^T,
\]
where $\bmP$ is a permutation matrix, $\bmL$ is unit lower triangular,
and $\bmB$ is block diagonal with blocks of size $1 \times 1$
and $2 \times 2$.

\noindent
2. $\bmA$ is {\bf positive definite}.
Now \packagename\ computes the {\bf sparse Cholesky factorization}
\[
\bmA =  \bmP\bmL(\bmP\bmL)^T,
\]
where $\bmP$ is a permutation matrix and $\bmL$ is lower triangular.
However, as \packagename\ is designed primarily for indefinite
systems, this may sometimes be slower than a dedicated Cholesky solver.

\noindent
There are options for iterative refinement and to
scale the matrix. In the latter case, the factorization of
the scaled matrix  $\overline{\bmA} = \bmS \bmA \bmS$ is computed,
where $\bmS$ is a diagonal scaling matrix.
The package returns bit-compatible results.

%%%%%%%%%%%%%%%%%%%%%% attributes %%%%%%%%%%%%%%%%%%%%%%%%

\galattributes
\galversions{\tt  \fullpackagename\_single, \fullpackagename\_double}.
\galuses {\tt \libraryname\_SMT}.
\galdate August 2026.
\galorigin J.\ Fowkes and N.\ I.\ M.\ Gould,
Rutherford Appleton Laboratory, and A.\ Montoison, AMD Inc., 
translated and extended as a fork of the open-source {\tt SSIDS} 
package, {\tt https://github.com/ralna/spral},
by J.\ Hogg and J.\ A.\ Scott.
\gallanguage Fortran~2003. 

%%%%%%%%%%%%%%%%%%%%%% HOW TO USE %%%%%%%%%%%%%%%%%%%%%%%%

\galhowto

\subsection{Calling sequences}

\input{versions}

\noindent
If it is required to use more than one of the modules at the same time, 
the derived types
{\tt SMT\_TYPE},
{\tt \packagename\_inform},
{\tt \packagename\_akeep},
{\tt \packagename\_fkeep}, and
{\tt \packagename\_control},
(Section~\ref{galtypes})
and the subroutines
{\tt \packagename\_analyse},
{\tt \packagename\_factor},
{\tt \packagename\_solve},
{\tt \packagename\_free},
(Section~\ref{galarguments})
{\tt \packagename\_enquire\_posdef},
{\tt \packagename\_enquire\_indef}, and
{\tt \packagename\_alter},
(Section~\ref{galfeatures})
must be renamed on one of the {\tt USE} statements.

\noindent There are five principal subroutines for user calls (see
\S\ref{galfeatures} for further features):

\begin{description}

\item[] The subroutine {\tt \packagename\_initialize} must be called to
specify the external solver to be used.
It may also be called to set default values
for solver-specific components of the control structure.
If non-default values are
wanted for any of the control components, the corresponding components
should be altered after the call to {\tt \packagename\_initialize}.

\item[] {\tt \packagename\_analyse} accepts the pattern of $\bmA$
 and optionally chooses pivots
 for Gaussian elimination using a selection criterion to preserve
 sparsity.  It subsequently constructs subsidiary information for
 actual factorization by {\tt \packagename\_factorize}. If the user provides
 the pivot sequence, only the necessary information for
 {\tt \packagename\_factorize} will be generated.

\item[] {\tt \packagename\_factorize} factorizes the matrix $\bmA$ using the
 information
 from a previous call to {\tt \packagename\_analyse}. The actual pivot sequence
 used may differ from that of {\tt \packagename\_analyse} if $\bmA$ is not
definite.

\item[] {\tt \packagename\_solve} uses the factors generated by
  {\tt \packagename\_factorize} to
  solve a system of equations with one ($\bmA \bmx = \bmr$)
  or several ($\bmA \bmX = \bmR$) right-hand sides.
  Iterative refinement may be used to improve a given solution
  or set of solutions.

\item[] {\tt \packagename\_terminate} deallocates the arrays held inside the
  structure for the factors. It should be called
  when all the systems involving its matrix have been solved
  or before another external solver is to be used.
\end{description}

%%%%%%%%%%%%%%%%%%%%%% matrix formats %%%%%%%%%%%%%%%%%%%%%%%%

\galmatrix
The matrix $\bmA$ may be stored in a variety of input formats.

\subsubsection{Sparse co-ordinate storage format}\label{coordinate}
Only the nonzero entries of the lower-triangular part of $\bmA$ are stored.
For the $l$-th entry of the lower-triangular portion of $\bmA$,
its row index $i$, column index $j$ and value $a_{ij}$
are stored in the $l$-th components of the integer arrays {\tt row},
{\tt col} and real array {\tt val}, respectively.
The order is unimportant, but the total number of entries
{\tt ne} is also required.

\subsubsection{Sparse row-wise storage format}\label{rowwise}
Again only the nonzero entries of the lower-triangular part are stored,
but this time they are ordered so that those in row $i$ appear directly
before those in row $i+1$. For the $i$-th row of $\bmA$, the $i$-th component
of an integer array {\tt ptr} holds the position of the first entry in this row,
while {\tt ptr} $(m+1)$ holds the total number of entries plus one.
The column indices $j$ and values $a_{ij}$ of the entries in the $i$-th row
are stored in components
$l =$ {\tt ptr}$(i)$, \ldots ,{\tt ptr} $(i+1)-1$ of the
integer array {\tt col}, and real array {\tt val}, respectively.

For sparse matrices, this scheme almost always requires less storage than
its predecessor. It is often known as CSR (compressed sparse-row) format.

\subsubsection{Dense storage format}\label{dense}
The matrix $\bmA$ is stored as a compact
dense matrix by rows, that is, the values of the entries of each row in turn are
stored in order within an appropriate real one-dimensional array.
Since $\bmA$ is symmetric, only the lower triangular part (that is the part
$a_{ij}$ for $1 \leq j \leq i \leq n$) need be held, and this part will be
stored by rows, that is component
$i \ast (i-1)/2 + j$ of the storage array {\tt val}
will hold the value $a_{ij}$ (and, by symmetry, $a_{ji}$)
for $1 \leq j \leq i \leq n$.

%\subsubsection{Diagonal storage format}\label{diagonal}
%If $\bmA$ is diagonal (i.e., $a_{ij} = 0$ for all $1 \leq i \neq j \leq n$)
%only the diagonals entries $a_{ii}$, $1 \leq i \leq n$,  need be stored,
%and the first $n$ components of the array {\tt A\%val} may be used for
%the purpose.

%%%%%%%%%%%%%%%%%%%%%%%%%%% kinds %%%%%%%%%%%%%%%%%%%%%%%%%%%

\input{kinds}

%%%%%%%%%%%%%%%%%%%%%% parallel usage %%%%%%%%%%%%%%%%%%%%%%%%

\input{omp}

%%%%%%%%%%%%%%%%%%%%%% derived types %%%%%%%%%%%%%%%%%%%%%%%%

\galtypes
Five derived data types are used by the package.

%%%%%%%%%%% problem type %%%%%%%%%%%

\subsubsection{The derived data type for holding the matrix}\label{typeprob}
The derived data type {\tt SMT\_type} is used to hold the matrix $\bmA$.
The components of {\tt SMT\_type} used are:

\begin{description}

\ittf{n} is a scalar variable of type \integer, that holds
the order $n$ of the matrix  $\bmA$.
\restriction {\tt n} $\geq$ {\tt 1}.

\itt{type} is an allocatable array of rank one and type default \character, that
indicates the storage scheme used. If the
sparse co-ordinate scheme (see \S\ref{coordinate}) is used
the first ten components of {\tt type} must contain the
string {\tt COORDINATE}.
For the sparse row-wise storage scheme (see \S\ref{rowwise}),
the first fourteen components of {\tt type} must contain the
string {\tt SPARSE\_BY\_ROWS}, and
for dense storage scheme (see \S\ref{dense})
the first five components of {\tt type} must contain the
string {\tt DENSE}.
%and for the diagonal storage scheme (see \S\ref{diagonal}),
%the first eight components of {\tt type} must contain the
%string {\tt DIAGONAL}.

For convenience, the procedure {\tt SMT\_put}
may be used to allocate sufficient space and insert the required keyword
into {\tt type}.
For example, if $\bmA$ is to be stored in the structure {\tt A}
of derived type {\tt SMT\_type} and we wish to use
the co-ordinate scheme, we may simply
%\vspace*{-2mm}
{\tt
\begin{verbatim}
        CALL SMT_put( A%type, 'COORDINATE', istat )
\end{verbatim}
}
%\vspace*{-4mm}
\noindent
See the documentation for the \galahad\ package {\tt SMT}
for further details on the use of {\tt SMT\_put}.

\itt{ne} is a scalar variable of type \integer, that
holds the number of entries in the {\bf lower triangular} part of $\bmA$
in the sparse co-ordinate storage scheme (see \S\ref{coordinate}).
It need not be set for any of the other three schemes.

\itt{val} is a rank-one allocatable array of type \realdp, that holds
the values of the entries of the {\bf lower triangular} part
of the matrix $\bmA$ for each of the
storage schemes discussed in \S\ref{galmatrix}.
Any duplicated entries that appear in the sparse
co-ordinate or row-wise schemes will be summed.

\itt{row} is a rank-one allocatable array of type \integer,
that holds the row indices of the {\bf lower triangular} part of $\bmA$
in the sparse co-ordinate storage
scheme (see \S\ref{coordinate}).
It need not be allocated for any of the other schemes.
Any entry whose row index lies out of the range $[$1,n$]$ will be ignored.

\itt{col} is a rank-one allocatable array variable of type \integer,
that holds the column indices of the {\bf lower triangular} part of
$\bmA$ in either the sparse co-ordinate
(see \S\ref{coordinate}), or the sparse row-wise
(see \S\ref{rowwise}) storage scheme.
It need not be allocated when the dense
%or diagonal storage schemes are used.
storage scheme is used.
Any entry whose column index lies out of the range $[$1,n$]$ will be ignored,
while the row and column indices of any entry from the
{\bf strict upper triangle} will implicitly be swapped.

\itt{ptr} is a rank-one allocatable array of size {\tt n+1} and type
\integer, that holds the starting position of
each row of the {\bf lower triangular} part of $\bmA$, as well
as the total number of entries plus one, in the sparse row-wise storage
scheme (see \S\ref{rowwise}). It need not be allocated for the
other schemes.

\end{description}
Although any of the above-mentioned matrix storage formats may be used
with each supported solver, {\tt MA27/SILS} and {\tt HSL\_MA57}
from HSL are most efficient if co-ordinate input is provided.

%%%%%%%%%%% control type %%%%%%%%%%%

\subsubsection{The derived data type for holding control
 parameters}\label{typecontrol}
The derived data type
{\tt \packagename\_control\_type}
is used to hold controlling data.
Default values specifically for the desired solver
may be obtained by calling
{\tt \packagename\_initialize}
(see \S\ref{subinit}),
%Default values are assigned,
while components may be changed at run time by calling
{\tt \packagename\_read\-\_specfile}
(see \S\ref{readspec}).
The components of
{\tt \packagename\_control\_type}
are:

\begin{description}

\itt{print\_level}
is a scalar variable of type \integer,
that gives the level of printing. Possible values are:
\begin{description}
\itt{<0} No printing.
\itt{0} Error and warning messages only.
\itt{1} As 0, plus basic diagnostic printing.
\itt{>1} As 1, plus some additional diagnostic printing.
\end{description}
The default is {\tt print\_level = 0}.

\itt{unit\_diagnostics}
is a scalar variable of type \integer,
that gives the Fortran unit number for diagnostics printing. 
Printing is suppressed if {\tt unit\_diagnostics} $<$0.
The default is {\tt unit\_diagnostics = 6}.

\itt{unit\_error}
is a scalar variable of type \integer,
that gives the Fortran unit number for printing
error messages. Printing is suppressed if {\tt unit\_error} $<$0.
The default is {\tt unit\_error = 6}.

\itt{unit\_warning}
is a scalar variable of type \integer,
that specifies the  Fortran unit number for printing of
warning messages. Printing is suppressed if {\tt unit\_warning} $<$0.
The default is {\tt unit\_warning = 6}.

\itt{ordering}
is a scalar variable of type \integer,
that gives the ordering method to use in analyse phase:
\begin{description}
\itt{0} User-supplied ordering is used ({\tt order} argument to
{\tt \packagename\_analyse} or {\tt \packagename\_analyse\_coord}).
\itt{1} MeTiS ordering with default settings.
\itt{2} Matching-based elimination ordering is computed (the
Hungarian algorithm is used to identify large
off-diagonal entries. A restricted MeTiS ordering is
then used that forces these on to the subdiagonal).
N.B.: This option should only be chosen for
indefinite systems. A scaling is also computed that may
be used in {\tt \packagename\_factor} (see \%scaling below).
\end{description}
The default is {\tt ordering = 1}.

\itt{nemin}
is a scalar variable of type \integer, that gives the
supernode amalgamation threshold. Two
neighbours in the elimination tree are merged if they both involve fewer
than nemin eliminations. The default is used if {\tt nemin $<$ 1}.
The default is {\tt nemin = 32}.

\itt{ignore\_numa}
is a scalar variable of type default \logical, that, if \true, specifies 
that all CPUs and GPUs are treated as belonging to a single NUMA region.
The default is {\tt ignore\_numa = \true}.

\itt{use\_gpu}
is a scalar variable of type default \logical, that, if \true, requests to
use an NVIDIA GPU, if present.
The default is {\tt use\_gpu = \true}.

\itt{min\_gpu\_work}
is a scalar variable of type \longinteger,
that gives the minimum number of flops
in subtree before scheduling on a GPU.
The default is {\tt min\_gpu\_work = 5.0E9}.

\itt{max\_load\_inbalance}
is a scalar variable of type \realdp, that specifies the maximum permissible 
load inbalance for leaf subtree allocations. Values less than 1.0 are treated
as 1.0.
The default is {\tt max\_load\_inbalance = 1.2}.

\itt{gpu\_perf\_coeff}
is a scalar variable of type \realdp, that holds the
GPU performance coefficient, that is how many
times faster a GPU is than CPU is at factoring a subtree.
The default is {\tt gpu\_perf\_coeff = 1.0}.

\itt{scaling}
is a scalar variable of type \integer,
that gives the scaling algorithm to use:
\begin{description}
\itt{$\leq$ 0} No scaling, if {\tt scale} is not present on call to
{\tt \packagename\_factor}, or user-supplied scaling (if
{\tt scale(:)} is present).
\itt{1} Compute using weighted bipartite matching via the
Hungarian Algorithm ({\tt MC64} algorithm).
\itt{2} Compute using a weighted bipartite matching via the
Auction Algorithm (may be lower quality than that
computed using the Hungarian Algorithm, but can be
considerably faster).
\itt{3} Use a matching-based ordering generated during the
analyse phase using {\tt control\%ordering=2}. The scaling
will be the same as that generated with
{\tt control\%scaling=1} if the matrix values have not
changed. This option will generate an error if a
matching-based ordering was not used during analysis.
\itt{$\geq$ 4} Compute using the norm-equilibration algorithm of Ruiz.
\end{description}
The default is {\tt scaling = 0}.

\itt{small\_subtree\_threshold}
is a scalar variable of type \longinteger, that gives the maximum number of
flops in a subtree treated as a single task (see Section~\ref{galmethod}).
The default is {\tt small\_subtree\_threshold = 4.OE6}.

\itt{cpu\_block\_size}
is a scalar variable of type \integer,
that gives the block size to use for parallelization of large nodes on 
CPU resources.
The default is {\tt cpu\_block\_size = 256}.

\itt{action}
is a scalar variable of type default \logical, that specifies whether to
continue factorization of singular matrix
on discovery of zero pivot if it is \true\ (a warning is issued), or to
abort if it is \false.
The default is {\tt action = \true}.

\itt{pivot\_method}
is a scalar variable of type \integer, that specifies the 
pivot method to be used on CPU, and can be one of:
\begin{description}
\itt{1} Aggressive a posteori pivoting. Cholesky-like
communication pattern is used, but a single failed pivot
requires restart of node factorization and potential
recalculation of all uneliminated entries.
\itt{2} Block a posteori pivoting. A failed pivot only requires
recalculation of entries within its own block column.
\itt{3} Threshold partial pivoting. Not parallel.
\end{description}
The default is {\tt pivot\_method = 2}.

\itt{small}
is a scalar variable of type \realdp, that holds the threshold below which an 
entry is treated as equivalent to {\tt 0.0}.
The default is {\tt small = 1E-20}.

\itt{u}
is a scalar variable of type \realdp, that specifies the
relative pivot threshold used in the symmetric
indefinite case. Values outside of the range $[0,0.5]$ are treated
as the closest value in that range.
The default is {\tt u = 0.01}.

\end{description}

%%%%%%%%%%% inform type %%%%%%%%%%%

\subsubsection{The derived data type for holding informational
 parameters}\label{typeinform}
The derived data type
{\tt \packagename\_inform\_type}
is used to hold parameters that give information about the progress and needs
of the algorithm. The components of
{\tt \packagename\_inform\_type}
are as follows:

\begin{description}

\itt{cpu\_flops}
is a scalar variable of type \longinteger,
that gives the number of flops performed on the CPUs.

\itt{cublas\_error}
is a scalar variable of type \integer,
that gives 

CUBLAS error code in the event of a CUBLAS error
(0 otherwise).

\itt{cuda\_error}
is a scalar variable of type \integer,
that gives the CUDA error code in the event of a CUDA error
(0 otherwise). Note that due to asynchronous execution, CUDA errors may
not be reported by the call that caused them.

\itt{flag}
is a scalar variable of type \integer,
that gives the exit status of the algorithm (see table below).

\itt{gpu\_flops}
is a scalar variable of type \longinteger,
number of flops performed on the GPUs.

\itt{matrix\_dup}
is a scalar variable of type \integer,
that gives the number of duplicate entries encountered (if
{\tt \packagename\_analyse} called with {\tt check=\true}, or any call to
{\tt \packagename\_analyse\_coord}).

\itt{matrix\_missing\_diag}
is a scalar variable of type \integer,
that gives the number of diagonal entries without an
explicit value (if {\tt \packagename\_analyse} called with {\tt check=\true}, or
any call to {\tt \packagename\_analyse\_coord}).

\itt{matrix\_outrange}
is a scalar variable of type \integer,
that gives the number of out-of-range entries encountered (if
{\tt \packagename\_analyse} called with {\tt check=\true}, or any call to
{\tt \packagename\_analyse\_coord}).

\itt{matrix\_rank}
is a scalar variable of type \integer,
that gives the (estimated) rank (structural after the analyse phase,
numerical after the factorization phase).

\itt{maxdepth}
is a scalar variable of type \integer,
that gives the maximum depth of the assembly tree.

\itt{maxfront}
is a scalar variable of type \integer,
that gives the maximum front size (without pivoting after the analyse
phase, with pivoting after the factorization phase).

\itt{maxsupernode}
is a scalar variable of type \integer,
that gives the maximum supernode size (without pivoting after the
analyse phase, with pivoting after the  factorization phase).

\itt{num\_delay}
is a scalar variable of type \integer,
that gives the number of delayed pivots. That is, the total
number of fully-summed variables that were passed to the father node
because of stability considerations. If a variable is passed further
up the tree, it will be counted again.

\itt{num\_factor}
is a scalar variable of type \longinteger,
that gives the number of entries in $L$ (without pivoting
after the analyse phase, and with pivoting after the factorization phase).

\itt{num\_flops}
is a scalar variable of type \longinteger,
that gives the number of floating-point operations for Cholesky
factorization (an indefinite factorization needs slightly more)
without pivoting after the analyse phase, with and with pivoting after 
the factorization phase.

\itt{num\_neg}
is a scalar variable of type \integer,
that gives the number of negative eigenvalues of the matrix $D$
after the factorization phase.

\itt{num\_sup}
is a scalar variable of type \integer,
that gives the number of supernodes in assembly tree.

\itt{num\_two}
is a scalar variable of type \integer,
that gives the number of $2 \times 2$ pivots used by the
factorization (i.e. in the matrix $B$).

\itt{stat}
is a scalar variable of type \integer,
that gives the Fortran allocation status parameter in event of allocation
error (0 otherwise).

\itt{status}
is a scalar variable of type \integer,
that gives the exit status of the algorithm.
A successful call is indicated by {\tt status = 0}.
See \S\ref{galerrors} for details of other values.

\end{description}

%%%%%%%%%%% data type %%%%%%%%%%%

\subsubsection{The derived data type for holding problem data}\label{typedata}
The derived data type
{\tt \packagename\_data\_type}
is used to hold all the data for a particular problem,
or sequences of problems with the same structure, between calls to
{\tt \packagename} procedures.
%This data should be preserved, untouched, from the initial call to
%{\tt \packagename\_initialize}
%to the final call to
%{\tt \packagename\_terminate}.
All components are private.

%%%%%%%%%%% data type %%%%%%%%%%%

\subsubsection{The derived data type for low-level problem data}
\label{typeakeep}
\label{typefkeep}
The derived data types
{\tt \packagename\_akeep\_type} and
{\tt \packagename\_akeep\_type}
may be used to hold the data for the symbolic and factorization phases 
for a particular problem,
or sequences of problems with the same structure, between calls to
low-level {\tt \packagename} procedures.
%This data should be preserved, untouched, from the initial call to
%{\tt \packagename\_initialize}
%to the final call to
%{\tt \packagename\_terminate}.
Once again, all components are private.

%%%%%%%%%%%%%%%%%%%%%% argument lists %%%%%%%%%%%%%%%%%%%%%%%%

\galarguments
We use square brackets {\tt [ ]} to indicate \optional\ arguments.

%%%%%% initialization subroutine %%%%%%

\subsubsection{The initialization subroutine}\label{subinit}
The initialization subroutine must be called for each solver used
to initialize data and solver-specific control parameters.

\hskip0.5in
{\tt CALL \packagename\_initialize( data, control, inform[, check] )}
\begin{description}

\itt{data} is a scalar \intentout argument of type
{\tt \packagename\_data\_type}
(see \S\ref{typedata}). It is used to hold data about the problem being
solved.

\itt{control} is a scalar \intentout argument of type
{\tt \packagename\_control\_type}
(see \S\ref{typecontrol}).
On exit, {\tt control} contains solver-specific default values for the
components as described in \S\ref{typecontrol}.
These values should only be changed after calling
{\tt \packagename\_initialize}.

\itt{inform} is a scalar \intentout argument of type
{\tt \packagename\_inform\_type}
(see \S\ref{typeinform}).
A successful call is indicated when the  component {\tt status} has the value 0.
For other return values of {\tt status}, see \S\ref{galerrors}.

\itt{check} is an \optional\ scalar \logical\ \intentin\ argument 
that if \present\ and set \true\ will check to see if the requested solver 
is available, and if not will replace this by a suitable equivalent; the 
equivalent will be recorded in {\tt control\%solver}. 
In addition, the output default values of
{\tt control\%ordering} and {\tt control\%scaling} may be adjusted to ensure
that these options are available.
No checks will be performed if {\tt check} is not \present\ or if it is set
to \false\ .
\end{description}

%%%%%%%%% analysis subroutine %%%%%%

\subsubsection{The sparsity pattern analysis subroutine}
The sparsity pattern of $\bmA$ may be analysed as follows:

\hskip0.5in
{\tt CALL \packagename\_analyse( matrix, data, control, inform[, PERM] )}

\begin{description}
\itt{matrix} is scalar \intentin\ argument of type {\tt SMT\_type}
that is used to specify $\bmA$.
The user must set all of the relevant components of {\tt matrix} according
to the storage scheme desired (see \S\ref{typeprob}) except
{\tt matrix\%val}. Incorrectly-set components will result in errors
flagged in {\tt inform\%status}, see \S\ref{galerrors}.

\itt{data} is a scalar \intentinout\ argument of type
{\tt \packagename\_data\_type}
(see \S\ref{typedata}). It is used to hold the factors and other
data about the problem being solved.
It must have been initialized by a call to
{\tt \packagename\_ini\-tialize}.

\itt{control} is scalar \intentin\ argument of type
{\tt \packagename\_control\_type}. Its components control the action
of the analysis phase, as explained in
\S\ref{typecontrol}.

\itt{inform} is a scalar \intentinout\ argument of type
{\tt \packagename\_inform\_type}
(see \S\ref{typeinform}).
A successful call is indicated when the  component {\tt status} has the value 0.
For other return values of {\tt status}, see \S\ref{galerrors}.

\itt{PERM} is an \optional\ \intentin\ rank-one \integer\ argument of
length {\tt matrix\%n} that may be used to provide a permutation/pivot
order.
If present, {\tt PERM(i)}, {\tt i} $= 1, \ldots,$ {\tt matrix\%n}, should be set
to the desired position of the input row $i$ in the permuted matrix.

\end{description}

%%%%%%%%% factorization subroutine %%%%%%

\subsubsection{The numerical factorization subroutine}
Once it has been analysed, the matrix $\bmA$ may be factorized as follows:

\hskip0.5in
{\tt CALL \packagename\_factorize(matrix, data, control, inform )}

\begin{description}

\itt{matrix} is scalar \intentin\ argument of type {\tt SMT\_type}
that is used to specify $\bmA$.
The user must set all of the relevant components of {\tt matrix} according
to the storage scheme desired (see \S\ref{typeprob}). Those components
set for {\tt \packagename\_analyse} must not have been altered in the interim.

\itt{data} is a scalar \intentinout\ argument of type
{\tt \packagename\_data\_type}
(see \S\ref{typedata}). It is used to hold the factors and other
data about the problem being solved.
It must have been initialized by a call to
{\tt \packagename\_ini\-tialize} and not altered by the user in the interim.

\itt{control} is scalar \intentin\ argument of type
{\tt \packagename\_control\_type}. Its components control the action
of the factorization phase, as explained in
\S\ref{typecontrol}.

\itt{inform} is a scalar \intentinout\ argument of type
{\tt \packagename\_inform\_type}
(see \S\ref{typeinform}).
A successful call is indicated when the  component {\tt status} has the value 0.
For other return values of {\tt status}, see \S\ref{galerrors}.

\end{description}

%%%%%%%%% solve subroutine %%%%%%

\subsubsection{The solution subroutine}
\label{solve}
Given the factorization, a set of equations may be solved as follows:

\hskip0.5in
{\tt CALL \packagename\_solve( matrix, X, data, control, inform )}

\begin{description}

\itt{matrix} is scalar \intentin\ argument of type {\tt SMT\_type}
that is used to specify $\bmA$.
The user must set all of the relevant components of {\tt matrix} according
to the storage scheme desired (see \S\ref{typeprob}). Those components
set for {\tt \packagename\_factorize} must not have been altered in the interim.

\ittf{X} is an \intentinout\ assumed-shape array argument of rank 1 or 2
and of type \realdp.  On entry, {\tt X} must be set
to the vector $\bmr$ or the matrix $\bmR$ and on successful return it holds
the solution $\bmx$ or $\bmX$. For the single right-hand side case, the
{\tt i}-th component of $\bmr$ and the resulting
{\tt i}-th component of the solution $\bmx$
occupy the {\tt i}-th component of {\tt X}. When there are multiple
right-hand sides, the
{\tt i}-th component of the $j$-th right-hand side $\bmr_j$
and the resulting solution $\bmx_j$ occupy the
{\tt i,j}-th component of {\tt X}.

\itt{data} is a scalar \intentinout\ argument of type
{\tt \packagename\_data\_type}
(see \S\ref{typedata}). It is used to hold the factors and other
data about the problem being solved.
It must have been initialized by a call to
{\tt \packagename\_ini\-tialize} and not altered by the user in the interim.

\itt{control} is scalar \intentin\ argument of type
{\tt \packagename\_control\_type}. Its components control the action
of the solve phase, as explained in
\S\ref{typecontrol}.

\itt{inform} is a scalar \intentinout\ argument of type
{\tt \packagename\_inform\_type}
(see \S\ref{typeinform}).
A successful call is indicated when the  component {\tt status} has the value 0.
For other return values of {\tt status}, see \S\ref{galerrors}.

\end{description}

%%%%%%% termination subroutine %%%%%%

\subsubsection{The termination subroutine}
All previously allocated internal arrays are deallocated and OpenMP locks
destroyed as follows:

\hskip0.5in
{\tt CALL \packagename\_terminate( data, control, inform )}

\begin{description}

\itt{data} is a scalar \intentinout\ argument of type
{\tt \packagename\_data\_type}
(see \S\ref{typedata}). It is used to hold the factors and other
data about the problem being solved.
It must have been initialized by a call to
{\tt \packagename\_ini\-tialize} and not altered by the user in the interim.
On exit, its allocatable array components will have been deallocated.

\itt{control} is scalar \intentin\ argument of type
{\tt \packagename\_control\_type}. Its components control the action
of the termination phase, as explained in
\S\ref{typecontrol}.

\itt{inform} is a scalar \intentinout\ argument of type
{\tt \packagename\_inform\_type}
(see \S\ref{typeinform}).
A successful call is indicated when the  component {\tt status} has the value 0.
For other return values of {\tt status}, see \S\ref{galerrors}.

\end{description}

%%%%%%%%%%%%%%%%%%%%%% Warning and error messages %%%%%%%%%%%%%%%%%%%%%%%%

\galerrors
A negative value of {\tt inform\%status} on exit from the subroutines
indicates that an error has occurred. No further calls should be made
until the error has been corrected. Possible values are:

\begin{description}

\itt{-1}
Error in sequence of calls (may be caused by failure of a preceding call).

\itt{-2}
{\tt n}$<$0 or {\tt ne}$<$1.

\itt{-3}
Error in {\tt ptr}.

\itt{-4} CSC format: All variable indices in one or more columns are
out-of-range.
Coordinate format: All entries are out-of-range.

\itt{-5} 
Matrix is singular and {\tt control\%action=\false}.

\itt{-6}
Matrix found not to be positive definite but {\tt posdef = \true}.

\itt{-7}
{\tt ptr} and/or {\tt row} is not present, but required as
{\tt \packagename\_analyse} was called with {\tt check=\false}.

\itt{-8}
{\tt control\%ordering} out of range, or {\tt control\%ordering=0} and
order parameter not provided or not a valid permutation.

\itt{-9}
{\tt control\%ordering=-2} but {\tt val} was not supplied.

\itt{-10}
ldx$<$n or nrhs$<$1.

\itt{-11}
job is out-of-range.

\itt{-13}
Called {\tt \packagename\_enquire\_posdef} on indefinite
factorization.

\itt{-14}
Called {\tt \packagename\_enquire\_indef} on positive-definite
factorization.

\itt{-15}
{\tt control\%scaling=3} but a matching-based ordering was not
performed during analyse phase.

\itt{-50}
Allocation error. If available, the stat parameter is
returned in inform\%stat.

\itt{-51}
CUDA error. The CUDA error return value is returned in
inform\%cuda\_error.

\itt{-52}
CUBLAS error. The CUBLAS error return value is returned in
inform\%cublas\_error.

\itt{-53}
OpenMP cancellation is disabled. Please set the environment
variable {\tt OMP\_CANCELLATION = \true}.

\end{description} 

A positive flag value is associated with a warning message.
Possible positive values are: 
 
\begin{description} 
\itt{+1}
Out-of-range variable indices found and ignored in input
data. inform\%matrix\_outrange is set to the number of such
entries.

\itt{+2}
Duplicate entries found and summed in input data.
inform\%matrix\_dup is set to the number of such entries.

\itt{+3}
Combination of +1 and +2.

\itt{+4}
One or more diagonal entries of $A$ are missing.

\itt{+5}
Combination of +4 and +1 or +2.

\itt{+6}
Matrix is found be (structurally) singular during analyse
phase. This will overwrite any of the above warning flags.

\itt{+7}
Matrix is found to be singular during the factorization phase.

\itt{+8}
Matching-based scaling found as side-effect of
matching-based ordering ignored
(consider setting {\tt control\%scaling=3}).

\itt{+50}
OpenMP processor binding is disabled. Consider setting
the environment variable {\tt OMP\_PROC\_BIND=\true} (this may
affect performance on NUMA systems).

\end{description}

%%%%%%%%%%%%%%%%%%%%%% Further control features %%%%%%%%%%%%%%%%%%%%%%%%

\galfeatures
\noindent In this section, we describe features for enquiring about and
manipulating the parts of the factorization constructed. These features
will not be needed by a user who wants simply to solve systems of
equations with matrix $\bmA$.

\packagename\ produces an $\bmL \bmB \bmL^T$ factorization
of $\bmA$ or a perturbation thereof, where $\bmL$ is a permuted lower
triangular matrix and $\bmB$ is a block diagonal matrix with blocks of order
1 and 2. It is convenient to write this factorization in the form
\disp{\bmA + \bmE = \bmP \bmL \bmB \bmL^T \bmP^T,}
where $\bmP$ is a permutation matrix and $\bmE$ is any diagonal perturbation
introduced. The following subroutines are provided:
\begin{description}
\ittf{} {\tt \packagename\_enquire} returns any of $\bmP$, $\bmB$, $\bmE$,
or the pivot permutations implicit in $\bmB$.

\ittf{} {\tt \packagename\_alter\_d} alters $\bmB$. Note that this means that
we no longer have a factorization of the given matrix $\bmA$.

\ittf{} {\tt \packagename\_part\_solve} solves one of the systems of equations
$\bmP \bmL \bmx =  \bmr$,
$\bmB \bmx =  \bmr$,
$\bmL^T \bmP^T \bmx =  \bmr$, or when $\bmA$ is positive definite,
$\bmP \bmL \sqrt{\bmB}\bmx =  \bmr$,
for one or more right-hand sides.

\end{description}

\subsubsection{To return $\bmP$ or $\bmB$ or both}
\label{galenquire}

\hskip0.5in
{\tt CALL \packagename\_enquire( data, inform[, PERM][, PIVOTS][, D][, PERTURBATION] )}

\begin{description}

\itt{data} is a scalar \intentinout\ argument of type
{\tt \packagename\_data\_type}
(see \S\ref{typedata}). It is used to hold the factors and other
data about the problem being solved.
It must have been initialized by a call to
{\tt \packagename\_ini\-tialize} and not altered by the user in the interim.

\itt{inform} is a scalar \intentinout\ argument of type
{\tt \packagename\_inform\_type}
(see \S\ref{typeinform}).
A successful call is indicated when the  component {\tt status} has the value 0.
For other return values of {\tt status}, see \S\ref{galerrors}.

\itt{PERM} is an \optional\ rank-one \integer\ array argument
of \intentout\ and length $n$.
If present, {\tt PERM} will be set to the pivot permutation
selected by {\tt \packagename\_analyse}.

\itt{PIVOTS} is an \optional\ rank-one \integer\ array argument
of \intentout\  and length $n$.
If present, the index of pivot $i$ will be placed in
{\tt PIVOTS} ($i$), $i = 1, \ldots, n$, with its sign negative if it is
the index of a 2 x 2 block.

\ittf{D} is an \optional\ rank-two \realdp\ array argument
of \intentout\, and shape {\tt (2, $n$)}.
If present, the diagonal entries of $\bmB^{-1} $ will
be placed in {\tt D( 1, i )}, $i = 1, \ldots n$ and the off-diagonal
entries of $\bmB^{-1} $ will be placed in {\tt D( 2, i )},
$i = 1, \ldots n-1$.

\itt{PERTURBATION} is an \optional\ rank-one \realdp\ array argument
of \intentout\ and length $n$.
If present, {\tt PERTURBATION} will be set to a vector of diagonal perturbations
chosen by {\tt \packagename\_factorize}. This array can only be nonzero if
{\tt \packagename\_factorize} was last called with {\tt control\%pivoting} = 4
or if the solver chose its own pivot sequence ({\tt control\%pivoting} $\leq 0$).

\end{description}

\subsubsection{To alter $\bmB$}

\hskip0.5in
{\tt CALL \packagename\_alter\_d( data, D, inform )}

\begin{description}

\itt{data} is a scalar \intentinout argument of type
{\tt \packagename\_data\_type}
(see \S\ref{typedata}). It is used to hold the factors and other
data about the problem being solved.
It must have been initialized by a call to
{\tt \packagename\_ini\-tialize} and not altered by the user in the interim.

\ittf{D} is an \intentinout\ \realdp\ array argument of shape ($2$, $n$).
The diagonal entries of $\bmB^{-1} $ will
be altered to {\tt D( 1, i )}, $i = 1, \ldots n$ and the off-diagonal
entries of $\bmB^{-1} $ will be altered to {\tt D( 2, i )},
$i = 1, \ldots n-1$.

\itt{inform} is a scalar \intentinout argument of type
{\tt \packagename\_inform\_type}
(see \S\ref{typeinform}).
A successful call is indicated when the  component {\tt status} has the value 0.
For other return values of {\tt status}, see \S\ref{galerrors}.

\end{description}

\subsubsection{To perform a partial solution}

\hskip0.5in
{\tt CALL \packagename\_part\_solve( part, X, data, control, inform )}

\begin{description}
\itt{part} is scalar, of \intentin\, and of type \character. It
must have one of the values
\begin{description}
\itt{L} for solving $\bmP \bmL \bmx =  \bmr$ or
                     $\bmP \bmL \bmX =  \bmR$,
\itt{B} for solving $\bmB \bmx =  \bmr$ or
                  $\bmB \bmX =  \bmR$, or
\itt{U} for solving $\bmL^T \bmP^T \bmx =  \bmr$ or
                  $\bmL^T \bmP^T \bmX = \bmR$.
\itt{S} for solving $\bmP \bmL \sqrt{\bmB} \bmx =  \bmr$ or
                  $\bmP \bmL \sqrt{\bmB} \bmX =  \bmR$
where $\sqrt{\bmB}$ is the diagonal matrix whose entries are the
square roots of those of $\bmB$ when $\bmA$ is positive definite.
\end{description}

\itt{data} is a scalar \intentinout argument of type
{\tt \packagename\_data\_type}
(see \S\ref{typedata}). It is used to hold the factors and other
data about the problem being solved.
It must have been initialized by a call to
{\tt \packagename\_ini\-tialize} and not altered by the user in the interim.

\ittf{X} is an \intentinout\ \realdp\ assumed-shape array argument of 
rank 1 or 2. On entry, {\tt X} must be set
to the vector $\bmr$ or the matrix $\bmR$ and on successful return it holds
the solution $\bmx$ or $\bmX$. For the single right-hand side case, the
{\tt i}-th component of $\bmr$ and the resulting
{\tt i}-th component of the solution $\bmx$
occupy the {\tt i}-th component of {\tt X}. When there are multiple
right-hand sides, the
{\tt i}-th component of the $j$-th right-hand side $\bmr_j$
and the resulting solution $\bmx_j$ occupy the
{\tt i,j}-th component of {\tt X}.

\itt{control} is scalar \intentin\ argument of type
{\tt \packagename\_control\_type}. Its components control the action
of the solve phase, as explained in
\S\ref{typecontrol}.

\itt{inform} is a scalar \intentinout\ argument of type
{\tt \packagename\_inform\_type}
(see \S\ref{typeinform}).
A successful call is indicated when the  component {\tt status} has the value 0.
For other return values of {\tt status}, see \S\ref{galerrors}.

\end{description}

%%%%%%%%%%%%%%%%%%%%%%%%%%% SOPHISTICATED %%%%%%%%%%%%%%%%%%%%%%%%%%%%%%%%%%

\section{Sophisticated usage}\label{sophisticated}
For convenience, \packagename\ also provides lower-level interfaces to its core
subroutines, more akin to those in provided by {\tt SSIDS}.

\subsection{Usage overview}
\label{usage-overview}

Solving $\bmA\bmX=\bmB$ using \packagename\ is a four stage process.
\begin{enumerate}
\item
Call {\tt \packagename\_analyse} to perform a symbolic factorization, stored
in {\tt akeep}.

\item
Call {\tt \packagename\_factor} to perform a numeric factorization, stored in
{\tt fkeep}. More than one numeric factorization can refer to the same 
{\tt akeep}.

\item
Call {\tt \packagename\_solve} to perform a solve with the factors. More than
one solve can be performed with the same {\tt fkeep}.

\item
Once all desired solutions have been found, free memory with
{\tt \packagename\_free}.

\end{enumerate}

In addition, advanced users may use the following routines:
\begin{itemize}
\item
{\tt \packagename\_enquire\_posdef} and 
{\tt \packagename\_enquire\_indef} return
the (block) diagonal entries of the factors and the pivot sequence.

\item
{\tt \packagename\_alter} modifies the (block) diagonal entries of the factors.
\end{itemize}

Note:
If used with bit-compatible BLAS and compiler options, this routine will
return bit compatible results. That is, consecutive runs with the same data
on the same machine produces exactly the same solution.


\subsection{Basic subroutines}
\label{basic-subroutines}

Note:
For the most efficient use of the package, Sparse row-wise storage (CSR) 
format (see \S\ref{rowwise}) should be used, without checking.

%%%%%%%%% analysis subroutine %%%%%%

\subsubsection{The analysis subroutine}

\hskip0.3in
{\tt CALL \packagename\_analyse( check, n, PTR, COL, akeep, control, inform{[}, ORDER, VAL, topology{]})}


Perform the analyse (symbolic) phase of the factorization for a matrix
supplied in CSR format. The resulting symbolic factors
stored in {\tt akeep} should be passed unaltered in the subsequent calls to
{\tt \packagename\_factor}.
\begin{description}
\itt{check}
is a scalar \intentin\ variable of type default \logical, 
that, if set \true\, causes checks of the matrix data.
Out-of-range entries are dropped and duplicate entries are summed.

\itt{n} is a scalar \intentin\ variable of type default \integer,
the gives the number of columns of $\bmA$.

\itt{PTR} is a rank-one array \intentin\ variable of type \longinteger\
and length at least {\tt n+1} that gives the starting addresses
for each row of the lower triangle of $\bmA$ in CSR format 
(see \S\ref{rowwise}), while {\tt ptr(n+1)} gives the total number of nonzeros
the lower triangle of $\bmA$ plus one.

\itt{COL} is a rank-one array \intentin\ variable of type \integer\
and length at least {\tt PTR(n+1)-1} that give the 
column indices for the lower triangle of $\bmA$ in CSR format.

\itt{akeep} is a scalar \intentinout argument of type
{\tt \packagename\_control\_akeep}
(see \S\ref{typeakeep}) that is used to store the 
symbolic factorization of $\bmA$, to be passed unchanged to subsequent routines.

\itt{control} is scalar \intentin\ argument of type
{\tt \packagename\_control\_type}. Its components control the action
of the analysis phase, as described in \S\ref{typecontrol}.

\itt{inform} is a scalar \intentinout\ argument of type
{\tt \packagename\_inform\_type}
(see \S\ref{typeinform}).
A successful call is indicated when the  component {\tt status} has the value 0.
For other return values of {\tt status}, see \S\ref{galerrors}.

\itt{order} is an \optional\ rank-one array \intentinout\ variable of type 
\integer\ and length at least {\tt n} that gives the values of 
a user-supplied ordering (see controll\%ordering=0) if required. 
On return, the actual ordering used is provided (if present).

\itt{VAL} is an \optional\ rank-one array \intentin\ variable of type \realdp\
and length at least {\tt PTR(n+1)-1} that give the values of the
nonzero entries for the lower triangle of $\bmA$ in CSR format.
This is only used if a matching-based ordering is requested.

\itt{topology} is an \optional\ rank-one allocatable \intentin\ array argument 
of type {\tt hw\_numa\_region} (see the package {\tt GALAHAD\_HW}).
If present, this specifies the machine topology
to be exploited. The size of the machine is the number of independent
NUMA regions. Region {\tt i} will use {\tt topology(i)\%nproc} threads and is
associated with the GPUs in the array {\tt topology(i)\%gpus}, which may have
length 0. If not present, these parameters are auto-detected using
the hwloc library (if detected at compiled time) or the environment
variable {\tt OMP\_NUM\_THREADS} if the hwloc library is not available. See
the method section for details of how work is divided.

\end{description}
Note:
If a user-supplied ordering is used, it may be altered by this routine,
with the altered version returned in order(:). This version will be
equivalent to the original ordering, except that some supernodes may have
been amalgamated, a topographic ordering may have been applied to the
assembly tree and the order of columns within a supernode may have been
adjusted to improve cache locality.

Note:
A version {\tt \packagename\_analyse\_ptr32},
where {\tt ptr} is of kind default \integer, is also provided for
backwards compatibility.

%%%%%%%%% analysis coodinate subroutine %%%%%%

\subsubsection{The analysis subroutine for coordinate storage}

Just as in {\tt \packagename\_analyse}, but for $\bmA$ stored in
sparse co-ordinate storage format (see \S\ref{coordinate}).

\hskip0.1in
{\tt CALL \packagename\_analyse\_coord( n, ne, ROW, COL, akeep, options, inform{[}, ORDER, VAL, topology{]})}

\noindent
The additional parameters are:
\begin{description}
\itt{ne} is a scalar \intentin\ variable of type default \integer,
the gives the number of nonzeros in the lower triangle of $\bmA$.

\itt{ROW} is a rank-one array \intentin\ variable of type \integer\
and length at least {\tt ne} that give the 
row indices for the lower triangle of $\bmA$ in co-ordinate format.

\itt{COL} is a rank-one array \intentin\ variable of type \integer\
and length at least {\tt ne} that give the 
column indices for the lower triangle of $\bmA$ in co-ordinate format.

\end{description}

%%%%%%%%% factorization subroutine %%%%%%

\subsubsection{The factorization subroutine}

\hskip0.5in
{\tt CALL \packagename\_factor( posdef, VAL, akeep, fkeep, options, inform{[}, SCALE, PTR, ROW{]})}

\begin{description}
\itt{posdef}
is a scalar \intentin\ variable of type default \logical, 
that should be \true\ if the matrix is known to be positive-definite.

\itt{VAL} is an \optional\ rank-one array \intentin\ variable of type \realdp\
that give the values of the nonzero entries for the lower triangle of $\bmA$ 
in the same format as for the call to {\tt \packagename\_analyse} or 
{\tt \packagename\_analyse\_coord}.

\itt{akeep} is a scalar \intentin argument of type
{\tt \packagename\_control\_akeep}
(see \S\ref{typeakeep}) that contains the symbolic factorization of $\bmA$
as returned by the preceding call
to {\tt \packagename\_analyse} or {\tt \packagename\_analyse\_coord}.

\itt{fkeep} is a scalar \intentinout argument of type
{\tt \packagename\_control\_fkeep}
(see \S\ref{typefkeep}) that is used to store the 
numerical factorization of $\bmA$, to be passed
unchanged to subsequent routines.

\itt{control} is scalar \intentin\ argument of type
{\tt \packagename\_control\_type}. Its components control the action
of the factorization phase, as described in \S\ref{typecontrol}.

\itt{inform} is a scalar \intentinout\ argument of type
{\tt \packagename\_inform\_type} (see \S\ref{typeinform}).
A successful call is indicated when the  component {\tt status} has the value 0.
For other return values of {\tt status}, see \S\ref{galerrors}.

\itt{SCALE} is an \optional\ rank-one array \intentinout\ variable of type 
\realdp\ and length at least {\tt n} that gives the values of 
user-supplied diagonal scalings if control\%scaling=0. Th ecomponent
{\tt SCALE(i)} should contain the entry $s_{ii}$ of $\bmS$. 
On return, the actual scalings used are provided (if present).

\itt{PTR} is an \optional\ rank-one array \intentin\ variable of type 
\longinteger\ and length at least {\tt n+1} that is only required if 
{\tt akeep} was obtained by running {\tt \packagename\_analyse} with
{\tt check} =.true., in which case it must be unchanged since that call.

\itt{COL} is an \optional\ rank-one array \intentin\ variable of type 
\integer\ and length at least {\tt PTR(n+1)-1} that is only required if 
{\tt akeep} was obtained by running {\tt \packagename\_analyse} with
{\tt check} =.true., in which case it must be unchanged since that call.

\end{description}

Note:
A version {\tt \packagename\_factor\_ptr32},
where {\tt ptr} is of kind default \integer, is also provided for
backwards compatibility.

%%%%%%%%% solve subroutine %%%%%%

\subsubsection{The solution subroutine for a single right-hand side}

Given factors of $\bmA$, solve (for a single right-hand side) a relevant 
system of equations.

\hskip0.5in
{\tt CALL \packagename\_solve( X, akeep, fkeep, options, inform{[}, job{]})}

\begin{description}

\itt{job}
is an \optional, \intentin\ scalar variable of type \integer,
that specifies the equation to be solved. 
Recall $\bmA$ has been factorized as either:
\begin{itemize}
\item{} 
$\bmS\bmA\bmS = (\bmP\bmL)(\bmP\bmL)^T~$ (positive-definite case); or
\item{} 
$\bmS\bmA\bmS = (\bmP\bmL)\bmB(\bmP\bmL)^T$ (indefinite case).
\end{itemize}
Possible values of {\tt job} are:
\begin{description}
\itt{0} (or absent) $\bmA\bmx=\bmr$
\itt{1} $\bmP\bmL\bmx=\bmS\bmr$
\itt{2} $\bmB\bmx=\bmr$
\itt{3} $(\bmP\bmL)^T\bmS^{-1}\bmx=\bmr$
\itt{4} $\bmB(\bmP\bmL)^T\bmS^{-1}\bmx=\bmr$
\end{description}

\itt{X} is a rank-one array \intentinout\ variable of type \realdp\ and 
length at least {\tt n} that holds the right-hand side $\bmr$ on entry, 
and provides the  solution $\bmx$ on exit.

\itt{akeep} is a scalar \intentin argument of type
{\tt \packagename\_control\_akeep} (see \S\ref{typeakeep}) that contains the 
symbolic factorization of $\bmA$ as returned by the preceding call
to {\tt \packagename\_analyse} or {\tt \packagename\_analyse\_coord}.

\itt{fkeep} is a scalar \intentinout argument of type
{\tt \packagename\_control\_fkeep} (see \S\ref{typefkeep}) that contains the 
numeric factorization returned by preceding call to {\tt \packagename\_factor}.

\itt{control} is scalar \intentin\ argument of type
{\tt \packagename\_control\_type}. Its components control the action
of the solve phase, as described in \S\ref{typecontrol}.

\itt{inform} is a scalar \intentinout\ argument of type
{\tt \packagename\_inform\_type} (see \S\ref{typeinform}).
A successful call is indicated when the  component {\tt status} has the value 0.
For other return values of {\tt status}, see \S\ref{galerrors}.

\end{description}

%%%%%%%%% multiple solve subroutine %%%%%%

\subsubsection{The solution subroutine for multiple right-hand sides}

Given factors of $\bmA$, solve (for a multiple right-hand sides) relevant 
system of equations.

\hskip0.5in
{\tt CALL \packagename\_solve( nrhs, X, ldx, akeep, fkeep, options, inform{[}, job{]})}

\begin{description}

\itt{job}
is an \optional, \intentin\ scalar variable of type \integer,
that specifies the equations to be solved. Possible values are:

\begin{description}
\itt{0 (or absent)} $\bmA\bmX=\bmR$
\itt{1} $\bmP\bmL\bmX=\bmS\bmR$
\itt{2} $\bmB\bmX=\bmR$
\itt{3} $(\bmP\bmL)^T\bmS^{-1}\bmX=\bmR$
\itt{4} $\bmB(\bmP\bmL)^T\bmS^{-1}\bmX=\bmR$
\end{description}

\itt{nrhs} is an \intentin\ scalar variable of type \integer,
that specifies the number of right-hand sides.

\itt{X} is a rank-two array \intentinout\ variable of type \realdp\ and 
dimensions {\tt ldx} and at least {\tt n}, respectively,
that holds the right-hand sides $\bmR$ on entry, and provides the 
solutions $\bmX$ on exit.

\itt{ldx} is an \intentin\ scalar variable of type \integer,
that specifies the leading dimension of {\tt X}.

\itt{akeep} is a scalar \intentin argument of type
{\tt \packagename\_control\_akeep} (see \S\ref{typeakeep}) that contains the 
symbolic factorization of $\bmA$ as returned by the preceding call
to {\tt \packagename\_analyse} or {\tt \packagename\_analyse\_coord}.

\itt{fkeep} is a scalar \intentinout argument of type
{\tt \packagename\_control\_fkeep} (see \S\ref{typefkeep}) that contains the 
numeric factorization returned by preceding call to {\tt \packagename\_factor}.

\itt{control} is scalar \intentin\ argument of type
{\tt \packagename\_control\_type}. Its components control the action
of the solve phase, as described in \S\ref{typecontrol}.

\itt{inform} is a scalar \intentinout\ argument of type
{\tt \packagename\_inform\_type} (see \S\ref{typeinform}).
A successful call is indicated when the  component {\tt status} has the value 0.
For other return values of {\tt status}, see \S\ref{galerrors}.

\end{description}

%%%%%%%%% termination subroutine %%%%%%

\subsubsection{The termination subroutine}

Free memory and resources associated with {\tt akeep} and/or
{\tt fkeep}, at least one of which must be present.

\hskip0.5in
{\tt CALL \packagename\_free({ [}akeep, fkeep,{]} cuda\_error )}

\begin{description}

\itt{akeep} is a scalar \intentinout argument of type
{\tt \packagename\_control\_akeep} (see \S\ref{typeakeep}) that contains the 
symbolic factorization of $\bmA$ as returned by the preceding call
to {\tt \packagename\_analyse} or {\tt \packagename\_analyse\_coord}.

\itt{fkeep} is a scalar \intentinout argument of type
{\tt \packagename\_control\_fkeep} (see \S\ref{typefkeep}) that contains the 
numeric factorization returned by preceding call to {\tt \packagename\_factor}.

\itt{cuda\_error} is an \intentout\ scalar variable of type \integer,
\integer that is {\tt 0} on a successful call, or a CUDA error code following
a failure.

\end{description}

Warning:
{\tt \packagename\_free} must be called by the user. Merely deallocating
{\tt akeep} or {\tt fkeep}, or allowing them to go out of scope
will result in memory leaks due to underlying memory allocation performed
using C++ and CUDA mechanisms. {\tt akeep} should only be deallocated
after all associated numeric factorizations {\tt fkeep} have been
freed.

\subsection{Advanced subroutines} \label{advanced-subroutines}

%%%%%%%%% enquire posdef subroutine %%%%%%

\subsubsection{The enquire subroutine for positive-definite matrices}

\hskip0.5in
{\tt CALL \packagename\_enquire\_posdef(akeep,fkeep,options,inform,D)}

Return the diagonal entries of the Cholesky factor.
\begin{description}
\itt{akeep} is a scalar \intentin argument of type
{\tt \packagename\_control\_akeep} (see \S\ref{typeakeep}) that contains the 
symbolic factorization of $\bmA$ as returned by the preceding call
to {\tt \packagename\_analyse} or {\tt \packagename\_analyse\_coord}.

\itt{fkeep} is a scalar \intentin argument of type
{\tt \packagename\_control\_fkeep} (see \S\ref{typefkeep}) that contains the 
numeric factorization returned by preceding call to {\tt \packagename\_factor}.

\itt{control} is scalar \intentin\ argument of type
{\tt \packagename\_control\_type}. Its components control the action
of the enquire phase, as described in \S\ref{typecontrol}.

\itt{inform} is a scalar \intentinout\ argument of type
{\tt \packagename\_inform\_type} (see \S\ref{typeinform}).
A successful call is indicated when the  component {\tt status} has the value 0.
For other return values of {\tt status}, see \S\ref{galerrors}.

\itt{D} is a rank-one array \intentout\ variable of type \realdp\ and 
length at least {\tt n} that returns the diagonal of $L$; {\tt d(i)} 
stores the entry $l_{ii}$.

\end{description}

%%%%%%%%% enquire indef subroutine %%%%%%

\subsubsection{The enquire subroutine for indefinite matrices}


\hskip0.5in
{\tt CALL \packagename\_enquire\_indef( akeep, fkeep, options, inform{[}, PIV\_ORDER, B{]})}


Return the pivot order and/or values of $\bmB$ of the Symmetric Indefinite
Factorization.
\begin{description}
\itt{akeep} is a scalar \intentin argument of type
{\tt \packagename\_control\_akeep} (see \S\ref{typeakeep}) that contains the 
symbolic factorization of $\bmA$ as returned by the preceding call
to {\tt \packagename\_analyse} or {\tt \packagename\_analyse\_coord}.

\itt{fkeep} is a scalar \intentin argument of type
{\tt \packagename\_control\_fkeep} (see \S\ref{typefkeep}) that contains the 
numeric factorization returned by preceding call to {\tt \packagename\_factor}.

\itt{control} is scalar \intentin\ argument of type
{\tt \packagename\_control\_type}. Its components control the action
of the enquire phase, as described in \S\ref{typecontrol}.

\itt{inform} is a scalar \intentinout\ argument of type
{\tt \packagename\_inform\_type} (see \S\ref{typeinform}).
A successful call is indicated when the  component {\tt status} has the value 0.
For other return values of {\tt status}, see \S\ref{galerrors}.

\itt{PIV\_ORDER} is an \optional\ rank-one array \intentout\ variable of type 
\integer\ and length at least {\tt n} that returns the pivot order.
The value $|${\tt piv\_order(i)}$|$ gives the position of variable
$i$ in the pivot order. The sign will be positive if $i$ is a
$1\times1$ pivot, and negative if $i$ is
part of a $2 \times 2$ pivot.

\itt{B} is an \optional\ rank-two array \intentout\ variable of type \realdp\ 
and dimensions {\tt 2} and at least {\tt n}, respectively,
returns the $2\times2$ block diagonals of $\bmB$. The component
{\tt b(1,i)} stores $b_{ii}$ and {\tt b(2,i)} stores $b_{(i+1)i}$.

\end{description}

%%%%%%%%% alter subroutine %%%%%%

\subsubsection{The alter subroutine}


Alter the entries of the diagonal factor $\bmB$ for a symmetric indefinite
factorization. The pivot order remains the same.

\hskip0.5in
{\tt CALL \packagename\_alter( B, akeep, fkeep, options, inform)}

\begin{description}
\itt{B} is an \optional\ rank-two array \intentout\ variable of type \realdp\ 
and dimensions {\tt 2} and at least {\tt n}, respectively,
stores the new $2\times2$ block diagonals of $\bmB$. The component
{\tt b(1,i)} stores $b_{ii}$ and {\tt b(2,i)} stores $b_{(i+1)i}$.

\itt{akeep} is a scalar \intentin argument of type
{\tt \packagename\_control\_akeep} (see \S\ref{typeakeep}) that contains the 
symbolic factorization of $\bmA$ as returned by the preceding call
to {\tt \packagename\_analyse} or {\tt \packagename\_analyse\_coord}.

\itt{fkeep} is a scalar \intentinout argument of type
{\tt \packagename\_control\_fkeep} (see \S\ref{typefkeep}) that contains the 
numeric factorization returned by preceding call to {\tt \packagename\_factor}
and records the change to $\\bmB$.

\itt{control} is scalar \intentin\ argument of type
{\tt \packagename\_control\_type}. Its components control the action
of the enquire phase, as described in \S\ref{typecontrol}.

\itt{inform} is a scalar \intentinout\ argument of type
{\tt \packagename\_inform\_type} (see \S\ref{typeinform}).
A successful call is indicated when the  component {\tt status} has the value 0.
For other return values of {\tt status}, see \S\ref{galerrors}.

\end{description}

%%%%%%%%%%%%%%%%%%%%%% Further features %%%%%%%%%%%%%%%%%%%%%%%%

\galcontrolfeatures
\noindent In this section, we describe an alternative means of setting
control parameters, that is components of the variable {\tt control} of type
{\tt \packagename\_control\_type}
(see \S\ref{typecontrol}),
by reading an appropriate data specification file using the
subroutine {\tt \packagename\_read\_specfile}. This facility
is useful as it allows a user to change  {\tt \packagename} control parameters
without editing and recompiling programs that call {\tt \packagename}.

A specification file, or specfile, is a data file containing a number of
``specification commands''. Each command occurs on a separate line,
and comprises a ``keyword'',
that is a string (in a close-to-natural language) used to identify a
control parameter, and
an (optional) "value", which defines the value to be assigned to the given
control parameter. All keywords and values are case insensitive,
keywords may be preceded by one or more blanks but
values must not contain blanks, and
each value must be separated from its keyword by at least one blank.
Values must not contain more than 30 characters, and
each line of the specification file is limited to 80 characters,
including the blanks separating keyword and value.

The portion of the specification file used by
{\tt \packagename\_read\_specfile}
must start
with a ``{\tt BEGIN \packagename}'' command and end with an
``{\tt END}'' command.  The syntax of the specfile is thus defined as follows:
\begin{verbatim}
  ( .. lines ignored by SLS_read_specfile .. )
    BEGIN SLS
       keyword    value
       .......    .....
       keyword    value
    END
  ( .. lines ignored by SLS_read_specfile .. )
\end{verbatim}
where keyword and value are two strings separated by (at least) one blank.
The ``{\tt BEGIN \packagename}'' and ``{\tt END}'' delimiter command lines
may contain additional (trailing) strings so long as such strings are
separated by one or more blanks, so that lines such as
\begin{verbatim}
    BEGIN SLS SPECIFICATION
\end{verbatim}
and
\begin{verbatim}
    END SLS SPECIFICATION
\end{verbatim}
are acceptable. Furthermore,
between the
``{\tt BEGIN \packagename}'' and ``{\tt END}'' delimiters,
specification commands may occur in any order.  Blank lines and
lines whose first non-blank character is {\tt !} or {\tt *} are ignored.
The content
of a line after a {\tt !} or {\tt *} character is also
ignored (as is the {\tt !} or {\tt *}
character itself). This provides an easy way to ``comment out'' some
specification commands, or to comment specific values
of certain control parameters.

The value of a control parameter may be of three different types, namely
integer, character or real.
Integer and real values may be expressed in any relevant Fortran integer and
floating-point formats (respectively).
%Permitted values for logical
%parameters are "{\tt ON}", "{\tt TRUE}", "{\tt .TRUE.}", "{\tt T}",
%"{\tt YES}", "{\tt Y}", or "{\tt OFF}", "{\tt NO}",
%"{\tt N}", "{\tt FALSE}", "{\tt .FALSE.}" and "{\tt F}".
%Empty values are also allowed for
%logical control parameters, and are interpreted as "{\tt TRUE}".

The specification file must be open for
input when {\tt \packagename\_read\_specfile}
is called, and the associated unit number
passed to the routine in {\tt device} (see below).
Note that the corresponding
file is rewound, which makes it possible to combine the specifications
for more than one program/routine.  For the same reason, the file is not
closed by {\tt \packagename\_read\_specfile}.

\subsubsection{To read control parameters from a specification file}
\label{readspec}

Control parameters may be read from a file as follows:
\hskip0.5in

\def\baselinestretch{0.8}
{\tt
\begin{verbatim}
     CALL SLS_read_specfile( control, device )
\end{verbatim}
}
\def\baselinestretch{1.0}

\begin{description}
\itt{control} is a scalar \intentinout\ argument of type
{\tt \packagename\_control\_type}
(see \S\ref{typecontrol}).
Default values should have already been set, perhaps by calling
{\tt \packagename\_initialize}.
On exit, individual components of {\tt control} may have been changed
according to the commands found in the specfile. Specfile commands and
the component (see \S\ref{typecontrol}) of {\tt control}
that each affects are given in Table~\ref{specfile}.

\pctable{|l|l|l|}
\hline
  command & component of {\tt control} & value type \\
\hline
{\tt error-printout-device} & {\tt \%error} & {\tt integer} \\
{\tt warning-printout-device} & {\tt \%warning} & {\tt integer} \\
{\tt printout-device} & {\tt \%out} & {\tt integer} \\
{\tt statistics-printout-device} & {\tt \%statistics} & {\tt integer} \\
{\tt print-level} & {\tt \%print\_level} & {\tt integer} \\
{\tt print-level-solver} & {\tt \%print\_level\_solver} & {\tt integer} \\
{\tt architecture-bits} & {\tt \%bits} & {\tt integer} \\
{\tt block-size-for-kernel} & {\tt \%block\_size\_kernel} & {\tt integer} \\
{\tt block-size-for--elimination} & {\tt \%block\_size\_elimination} & {\tt integer} \\
{\tt blas-block-for-size-factorize} & {\tt \%blas\_block\_size\_factorize} & {\tt integer} \\
{\tt blas-block-size-for-solve} & {\tt \%blas\_block\_size\_solve} & {\tt integer} \\
{\tt node-amalgamation-tolerance} & {\tt \%node\_amalgamation} & {\tt integer} \\
{\tt initial-pool-size} & {\tt \%initial\_pool\_size} & {\tt integer} \\
{\tt minimum-real-factor-size} & {\tt \%min\_real\_factor\_size} & {\tt integer} \\
{\tt minimum-integer-factor-size} & {\tt \%min\_integer\_factor\_size} & {\tt integer} \\
{\tt maximum-real-factor-size} & {\tt \%max\_real\_factor\_size} & {\tt integer(long)} \\
{\tt maximum-integer-factor-size} & {\tt \%max\_integer\_factor\_size} & {\tt integer(long)} \\
{\tt maximum-in-core-store} & {\tt \%max\_in\_core\_store} & {\tt integer(long)} \\
{\tt pivot-control} & {\tt \%pivot\_control} & {\tt integer} \\
{\tt ordering} & {\tt \%ordering} & {\tt integer} \\
{\tt full-row-threshold} & {\tt \%full\_row\_threshold} & {\tt integer} \\
{\tt pivot-row-search-when-indefinite} & {\tt \%row\_search\_indefinite} & {\tt integer} \\
{\tt scaling} & {\tt \%scaling} & {\tt integer} \\
{\tt scale-maxit} & {\tt \%scale\_maxit} & {\tt integer} \\
{\tt scale-thresh} & {\tt \%scale\_thresh} & {\tt real} \\
{\tt max-iterative-refinements} & {\tt \%max\_iterative\_refinements} & {\tt integer} \\
{\tt array-increase-factor} & {\tt \%array\_increase\_factor} & {\tt real} \\
{\tt array-decrease-factor} & {\tt \%array\_decrease\_factor} & {\tt real} \\
{\tt relative-pivot-tolerance} & {\tt \%relative\_pivot\_tolerance} & {\tt real} \\
{\tt absolute-pivot-tolerance} & {\tt \%absolute\_pivot\_tolerance} & {\tt real} \\
{\tt zero-tolerance} & {\tt \%zero\_tolerance} & {\tt real} \\
{\tt static-pivot-tolerance} & {\tt \%static\_pivot\_tolerance} & {\tt real} \\
{\tt static-level-switch} & {\tt \%static\_level\_switch} & {\tt real} \\
{\tt consistency-tolerance} & {\tt \%consistency\_tolerance} & {\tt real} \\
{\tt acceptable-residual-relative} & {\tt \%acceptable\_residual\_relative} & {\tt real} \\
{\tt acceptable-residual-absolute} & {\tt \%acceptable\_residual\_absolute} & {\tt real} \\
{\tt out-of-core-directory} & {\tt \%out\_of\_core\_directory} & {\tt character} \\
{\tt out-of-core-integer-factor-file} & {\tt \%out\_of\_core\_integer\_factor\_file} & {\tt character} \\
{\tt out-of-core-real-factor-file} &  {\tt \%out\_of\_core\_real\_factor\_file} & {\tt character} \\
{\tt out-of-core-real-work-file} &  {\tt \%out\_of\_core\_real\_work\_file} & {\tt character} \\
{\tt out-of-core-indefinite-file} &  {\tt \%out\_of\_core\_indefinite\_file} & {\tt character} \\
{\tt output-line-prefix} & {\tt \%prefix} & {\tt character} \\
\hline

\ectable{\label{specfile}Specfile commands and associated
components of {\tt control}.}

\itt{device} is a scalar \intentin\ argument of type \integer,
that must be set to the unit number on which the specification file
has been opened. If {\tt device} is not open, {\tt control} will
not be altered and execution will continue, but an error message
will be printed on unit {\tt control\%error}.

\end{description}

%%%%%%%%%%%%%%%%%%%%%% GENERAL INFORMATION %%%%%%%%%%%%%%%%%%%%%%%%

\galgeneral

%\galcommon None.
\galworkspace Provided automatically by the module.
\galmodules {\tt GALAHAD\_CLOCK},
{\tt GALAHAD\_SYMBOLS},
{\tt GALAHAD\_SORT\_single/double},
{\tt GALAHAD\_SPACE\-\_single/double},
{\tt GALAHAD\_SPECFILE\_single/double},
{\tt GALAHAD\_STRING\_single/double},
{\tt GALAHAD\_SMT\-\_sin\-gle/double},
{\tt GALAHAD\_SILS\_single/double},
{\tt GALAHAD\_NODEND\_single/double}
and optionally
{\tt HSL\_MA57\_single/double},
{\tt HSL\_MA77\_single/\-double},
{\tt HSL\_MA86\_single/double},
{\tt HSL\_MA87\_single/double},
{\tt HSL\_MA97\_single/double},
{\tt HSL\_MC64\_single\-/double} and
{\tt HSL\_MC68\_single\-/double}.
\galroutines Optionally {\tt MC61}, {\tt MC77} and {\tt METIS}.
\galio Output is under control of the arguments
{\tt control\%error},
{\tt control\%warning},
{\tt control\%out}, \\
{\tt control\-\%statistics}
and {\tt control\%print\_level}.
\galrestrictions {\tt matrix\%n} $\geq$ {\tt 1},
{\tt matrix\%ne} $\geq$ {\tt 0} if
{\tt matrix\%type = 'COORDINATE'},
{\tt matrix\%type}
one of
{\tt 'COORDINATE'}, {\tt 'SPARSE\_BY\_ROWS'} or   {\tt 'DENSE'}.
% or {\tt 'DIAGONAL'}.
\galportability ISO Fortran~95 + TR 15581 or Fortran~2003 and optionally OpenMP.
The package is thread-safe.

%%%%%%%%%%%%%%%%%%%%%% METHOD %%%%%%%%%%%%%%%%%%%%%%%%

\galmethod

\subsection{Partition of work across available resources}
\label{partition-of-work-across-available-resources}
Once the ordering has been determined and the assembly tree determined in the
analyse phase, the tree is broken into a number of leaf subtrees rooted at a
single node, leaving a root subtree/forest consisting of all remaining nodes
above those. Each leaf subtree is pre-assigned to a particular 
NUMA region or GPU for the factorization phase. Details of the algorithm 
used for finding these subtrees and their assignment can be found in 
Duff, Hogg and Lopez (2020), \S~\ref{galreferences}.

The factorization phase has two steps. In the first, leaf subtrees are
factorized in parallel on their assigned resources. Once all leaf subtrees are
factored, the second phase begins where all CPU resources cooperate to factorize
the root subtree.

At present the solve phase is performed in serial.

\subsection{Data checking}
%\label{\detokenize{data-checking}
If {\tt check} is set to \true\ on the call to {\tt \packagename\_analyse} or if
{\tt \packagename\_analyse\_coord} is called, the user-supplied matrix data is
checked for errors. The cleaned integer matrix data (duplicates are
summed and out-of-range indices discarded) is stored in akeep. The use
of checking is optional on a call to {\tt \packagename\_analyse} as it incurs 
both time and memory overheads. However, it is recommended since the
behaviour of the other routines in the package is unpredictable if
duplicates and/or out-of-range variable indices are entered.

If the user has supplied an elimination order it is checked for errors.
Otherwise, an elimination order is generated by the package. The
elimination order is used to construct an assembly tree. On exit from
{\tt \packagename\_analyse} (and {\tt \packagename\_analyse\_coord}), 
{\tt order} is
set so that {\tt order(i)} holds the position of variable $i$ in the
elimination order. If an ordering was supplied by the user, this order may
differ, but will be equivalent in terms of fill-in.


\subsection{Factorization performed}
%\label{factorization-performed}}
The factorization performed depends on the value of {\tt posdef}.

If {\tt posdef=\true}, a Cholesky factorization is performed:
\[
\bmS\bmA\bmS = \bmP\bmL(\bmP\bmL)^T.
\]
Pivoting is not performed, so $\bmP$ is the permutation determined in the
analysis phase. $\bmS$ is a diagonal scaling matrix.

If {\tt posdef=\false}, a symmetric indefinite factorization is performed:
\[
\bmS\bmA\bmS = \bmP\bmL\bmB(\bmP\bmL)^T
\]
Pivoting is performed, so $\bmP$ may differ from the permutation determined
in the analysis phase, though it is kept as close as possible to minimize fill.
The exact pivoting algorithm varies depending on the particular kernel the
algorithm chooses to employ.

Full details of the algorithms used are given by
Hogg, Ovtchinnikov and Scott (2014) for GPUs and 
Duff, Hogg and Lopez (2020) for CPUs (see \S~\ref{galreferences}).

\subsection{Small Leaf Subtrees}
\label{small-leaf-subtrees}\label{small-leaf}
For subtrees allocated to run on the CPU, the factorization of small nodes near
the leaves of the tree can be amalgamated into a single parallel task (normally
each would be treated as its own OpenMP task to be scheduled). This can reduce
scheduling overheads, especially on small problems. If the total number of
operations for a subtree root at a given node is less than
{\tt control\%small\_subtree\_threshold}, that subtree is treated as a 
single task.

\galreferences
\vspace*{1mm}

\noindent
The detailed algorithm is described for GPUs in
\vspace*{1mm}

\noindent
J.\ D.\ Hogg, E.\ Ovtchinnikov and J.\ A.\ Scott (2014).
A sparse symmetric indefinite direct solver for GPU architectures.
ACM Transactions on Mathematical Software 42(1), Article 1, 25 pages,
\vspace*{1mm}

\noindent
and its adaptation for CPUs in

\noindent
I.\ S.\ Duff, J.\ D.\ Hogg and F.\ Lopez (2020).
A New Sparse $LDL^T$ Solver Using A Posteriori Threshold Pivoting,
SIAM Journal on Scientific Computing 42(2), C23-C42.

%%%%%%%%%%%%%%%%%%%%%% EXAMPLE %%%%%%%%%%%%%%%%%%%%%%%%

\galexample
We illustrate the use of the package on the solution of the
single set of equations
\disp{\mat{ccccc}{ 2 & 3 &   &   &   \\
                   3 &   & 4 &   & 6 \\
                     & 4 & 1 & 5 &   \\
                     &   & 5 &   &   \\
                     & 6 & &     & 1 } \bmx =
       \vect{ 8 \\ 45 \\ 31 \\ 15 \\ 17}}
(Note that this example does not illustrate all the facilities).



\end{document}




{\tt \small
\VerbatimInput{\packageexample}
}
\noindent
%with the following data
%{\tt \small
%\VerbatimInput{\packagedata}
%}
%\noindent
This produces the following output:
{\tt \small
\VerbatimInput{\packageresults}
}
\noindent
The same problem may be solved holding the data in
a sparse row-wise storage format by replacing the lines
{\tt \small
\begin{verbatim}
! allocate and set lower triangle of matrix in co-ordinate form
...
! problem setup complete
\end{verbatim}
}
\noindent
by
{\tt \small
\begin{verbatim}
! allocate and set lower triangle of matrix in spares row form
     CALL SMT_put( matrix%type, 'SPARSE_BY_ROWS', s )
     matrix%n = n 
     ALLOCATE( matrix%val( ne ), matrix%col( ne ), matrix%ptr( n + 1 ) )
     matrix%ptr = (/ 1, 2, 3, 5, 6, 8 /)
     matrix%col = (/ 1, 1, 2, 3, 3, 2, 5 /)
     matrix%val = (/ 2.0_wp, 3.0_wp, 4.0_wp, 1.0_wp, 5.0_wp, 6.0_wp, 1.0_wp /)
! problem setup complete
\end{verbatim}
}
\noindent
or using a dense storage format with the replacement lines
{\tt \small
\begin{verbatim}
! allocate and set lower triangle of matrix in dense form
     CALL SMT_put( matrix%type, 'DENSE', s )
     matrix%n = n 
     ALLOCATE( matrix%val( n * ( n + 1 ) / 2 ) )
     matrix%val = (/ 2.0_wp, 3.0_wp, 0.0_wp, 0.0_wp, 4.0_wp, 1.0_wp,           &
                     0.0_wp, 0.0_wp, 5.0_wp, 0.0_wp, 0.0_wp, 6.0_wp,           &
                     0.0_wp, 0.0_wp, 1.0_wp /)
! problem setup complete
\end{verbatim}
}
\noindent
respectively.


\end{document}
